\section{All-order local identities and finite differences}
\label{asj:sec:local}

Normalize the exponent vector by $\sum_j\alpha_j=1$, and set
\begin{equation}\label{asj:eq:Jnormalized}
 \calJ_r(\boldsymbol\alpha)
  =r![s^r]\calZ_P(s\boldsymbol\alpha),\qquad r\ge0.
\end{equation}
Let $F_{[q]}$ denote the homogeneous Taylor part of degree $q$ of a
holomorphic germ or polynomial $F$. Put
$H=\max_{p\in\mathcal P}H_p$, with $H=0$ if $\mathcal P$ is empty.

\begin{theorem}[Polynomiality of every normalized Laurent jet]
\label{asj:thm:polynomial}
On the affine hyperplane $\sum\alpha_j=1$,
\begin{align}
 \calJ_r(\boldsymbol\alpha)
  &=r!\calH_{P,[r]}(\boldsymbol\alpha)
      +r!\sum_{p,h}\kappa_{p,h}
                       \calR_{p,[r+h]}(\boldsymbol\alpha),\label{asj:eq:Jlocal}\\
 \calR_{p,[q]}(\boldsymbol\alpha)
  &=\frac{(-1)^q}{q!}[x^{-p}]P(x)
       \left(\sum_j\alpha_j\log(1+a_j/x)\right)^q.
 \label{asj:eq:Rhom}
\end{align}
Consequently $\deg\calJ_r\le r+H$. A local term of degree $r+h$ can
occur only when $r+h\le d+p$. Every term beyond degree $r$ is determined
by the finite principal parts of $\D$.
\end{theorem}
\begin{proof}
On the normalized ray $U=s$. Extracting the coefficient of $s^r$ in
\eqref{asj:eq:germ} gives \eqref{asj:eq:Jlocal}. The homogeneous part of degree
$q$ of $e^{-L_{\boldsymbol u}(z)}$ is
$(-1)^qL_{\boldsymbol u}(z)^q/q!$. Using \eqref{asj:eq:Rdef} and then
$z=x^{-1}$ proves \eqref{asj:eq:Rhom}. Every factor of the logarithmic
sum starts with $x^{-1}$, proving the sharper degree cutoff.
\end{proof}

For a tangent vector $\boldsymbol v$ with $\sum_jv_j=0$, write
\[
 \Delta_{\boldsymbol v}f(\boldsymbol\alpha)
     =f(\boldsymbol\alpha+\boldsymbol v)-f(\boldsymbol\alpha),
 \qquad V_{\boldsymbol v}(x)=\sum_jv_j\log(x+a_j).
\]
The $\log x$ terms cancel, so $V_{\boldsymbol v}(x)=O(x^{-1})$.

\begin{theorem}[Top local polarization]\label{asj:thm:top}
Suppose $H\ge1$ and let $q=r+H$. For any $q$ tangent vectors,
\begin{align}\label{asj:eq:top}
 &\Delta_{\boldsymbol v_1}\cdots\Delta_{\boldsymbol v_q}
          \calJ_r(\boldsymbol\alpha)\notag\\
 &\quad=(-1)^q r!\sum_{p:H_p=H}\kappa_{p,H}
           [x^{-p}]P(x)\prod_{i=1}^q V_{\boldsymbol v_i}(x).
\end{align}
The same expression equals the corresponding mixed directional
derivative. Every mixed difference of order $q+1$ vanishes.
\end{theorem}
\begin{proof}
The $q$ differentiations annihilate $\calH_{P,[r]}$ and every local
part of degree below $q$. Polarizing the remaining $q$th power in
\eqref{asj:eq:Rhom} supplies exactly $q!\prod_i V_{\boldsymbol v_i}$.
For a polynomial of degree at most $q$, a mixed $q$th finite difference
is its constant mixed $q$th derivative; integrate that derivative over
the unit $q$-cube to see the equality directly. Differences of the
next order vanish by the degree bound.
\end{proof}

The highest pole alone controls \eqref{asj:eq:top}; it does not control
all lower local differences. The following formula retains every layer.

\begin{theorem}[All finite local differences]\label{asj:thm:all-differences}
Let $q>r$, and put $L_{\boldsymbol\alpha}(x)=
\sum_j\alpha_j\log(1+a_j/x)$. Then
\begin{align}\label{asj:eq:all-differences}
 &\Delta_{\boldsymbol v_1}\cdots\Delta_{\boldsymbol v_q}
       \calJ_r(\boldsymbol\alpha)\notag\\
 &=r!\sum_{p,h}\kappa_{p,h}[x^{-p}]P(x)
      [t^{r+h}]e^{-tL_{\boldsymbol\alpha}(x)}
            \prod_{i=1}^q\bigl(e^{-tV_{\boldsymbol v_i}(x)}-1\bigr).
\end{align}
The right side is a finite algebraic computation. Terms with $q>r+h$
vanish automatically.
\end{theorem}
\begin{proof}
The differences kill the degree-$r$ holomorphic part. Apply a difference
to $e^{-tL_{\boldsymbol\alpha}}$: it multiplies the exponential by
$e^{-tV_{\boldsymbol v}}-1$. Repeating and extracting $t^{r+h}$ in
\eqref{asj:eq:Jlocal} gives the assertion. Each of the $q$ factors is
divisible by $t$, proving the last statement. Coefficients in $x^{-1}$
need only be retained through the largest integer $d+p$ occurring.
\end{proof}

\subsection{Normalization, scaling, and the meaning of locality}
For arbitrary weights $\boldsymbol w$ with $W=\sum_jw_j\ne0$,
\begin{equation}\label{asj:eq:scaling}
 r![s^r]\calZ_P(s\boldsymbol w)
   =W^r\calJ_r(\boldsymbol w/W).
\end{equation}
This remains true for Laurent series: substitute $t=Ws$ before extracting
the coefficient. The excluded hyperplane $W=0$ cannot be restored by
blindly taking a limit in these formulas; it lies in the polar divisor
and needs a separately specified restriction.

``Local'' here has a precise meaning: the expression depends on finitely
many principal coefficients of the base Dirichlet series and finitely many
terms of the large-$x$ logarithmic expansion. It is not a claim that the
individual jets have elementary evaluations. Finite-head changes affect
only $\calH_P$, so any identity obtained by more than $r$ tangent differences
is insensitive to finitely many altered spectral terms.
